\section{Введение}

\subsection{Актуальность}

В современных распределённых системах REST API является основным способом взаимодействия между сервисами. При обработке HTTP-запросов неизбежно возникают ситуации, когда запрос может быть отправлен повторно: из-за сетевых таймаутов, retry-политик, дублирования на уровне балансировщиков нагрузки или ошибок клиента. Без механизмов идемпотентности повторная обработка запроса приводит к критическим последствиям: дублированию платежей, двойному списанию средств, созданию дубликатов записей в базе данных~\cite{helland2012idempotence}.

Тестирование идемпотентности особенно важно в финансовых и e-commerce системах, где повторная обработка транзакции приводит к прямым финансовым потерям. Компании, такие как Stripe, реализуют паттерн Idempotency-Key как обязательный элемент API~\cite{stripe-idempotency}.

\subsection{Постановка проблемы}

Проблема заключается в необходимости обеспечения корректного поведения сервиса при повторных запросах. HTTP-спецификация~\cite{rfc9110} определяет, какие методы должны быть идемпотентными (GET, PUT, DELETE), а какие --- нет (POST). Однако на практике реализация идемпотентности требует дополнительных механизмов на уровне приложения, причём конкретный механизм существенно зависит от выбранного транспорта (REST, SOAP, gRPC, GraphQL) и от стиля организации контроллера в Spring.

Цель данной работы --- исследовать и продемонстрировать методы тестирования идемпотентности и повторяемости запросов на примере Spring Boot приложения с шестью кардинально различными типами организации контроллеров, охватывающими четыре транспортных протокола: HTTP/REST (в трёх стилях), SOAP, gRPC и GraphQL.

\subsection{Задачи}

\begin{enumerate}
    \item Формализовать понятие идемпотентности HTTP-запросов математически.
    \item Разработать сервис с шестью различными типами контроллеров, объединёнными единым механизмом идемпотентности:
    \begin{itemize}
        \item аннотационный \texttt{@RestController} (по образцу Spring Guides~\cite{spring-guides-rest});
        \item функциональные эндпоинты \texttt{RouterFunction} (по образцу Spring Framework~\cite{spring-functional-endpoints});
        \item автогенерируемые эндпоинты Spring Data REST~\cite{spring-data-rest-guide};
        \item SOAP-эндпоинт на Spring-WS~\cite{spring-ws} (контракт-первый, XSD $\to$ WSDL);
        \item gRPC-сервис на базе \texttt{grpc-spring-boot-starter}~\cite{grpc-java} и сгенерированных protobuf-стабов;
        \item GraphQL-контроллер на \texttt{spring-boot-starter-graphql}~\cite{spring-graphql} со схемой SDL.
    \end{itemize}
    \item Реализовать механизмы обеспечения идемпотентности (Idempotency-Key, optimistic locking, unique constraints) и продемонстрировать их переиспользование между транспортами: единый \texttt{IdempotencyService} обслуживает HTTP-фильтр (REST/SOAP/GraphQL) и gRPC ServerInterceptor.
    \item Сделать сам механизм \emph{транспорт-агностичным}: модернизировать существующий фильтр так, чтобы он корректно реплеил ответы любого формата (JSON, XML/SOAP, бинарный protobuf).
    \item Разработать набор интеграционных тестов с использованием Testcontainers~\cite{testcontainers}.
    \item Продемонстрировать поведение системы без механизмов идемпотентности через отдельный Spring-профиль, причём для каждого из шести транспортов.
\end{enumerate}
